% XeLaTeX can use any Mac OS X font. See the setromanfont command below.
% Input to XeLaTeX is full Unicode, so Unicode characters can be typed directly into the source.

% The next lines tell TeXShop to typeset with xelatex, and to open and save the source with Unicode encoding.

%!TEX TS-program = xelatex
%!TEX encoding = UTF-8 Unicode

\documentclass[12pt,twoside]{article}

\usepackage{geometry}                % See geometry.pdf to learn the layout options. There are lots.
\geometry{a4paper}                   % ... or a4paper or a5paper or ... 
%\geometry{landscape}                % Activate for for rotated page geometry
%\usepackage[parfill]{parskip}    % Activate to begin paragraphs with an empty line rather than an indent
\usepackage{amssymb}

% Will Robertson's fontspec.sty can be used to simplify font choices.
% To experiment, open /Applications/Font Book to examine the fonts provided on Mac OS X,
% and change "Hoefler Text" to any of these choices.

\usepackage{fontspec,xltxtra,xunicode}
\defaultfontfeatures{Mapping=tex-text}
\setromanfont[Mapping=tex-text]{Hoefler Text}
\setsansfont[Scale=MatchLowercase,Mapping=tex-text]{Optima}
\setmonofont[Scale=MatchLowercase]{Andale Mono}

%%%% Graphics
\usepackage{graphicx}
\usepackage{tikz}
\definecolor{unilublue}{RGB}{55,149,218}
\definecolor{sntred}{RGB}{219,46,27}
\definecolor{sntpurple}{RGB}{86,30,130}
%\definecolor{sntblue}{RGB}{50,130,207}
\definecolor{sntblue}{RGB}{55,149,218}

\pgfdeclareimage[width=30mm]{logo-snt}{logos/logo-snt}
\pgfdeclareimage[width=30mm]{logo-uni-lu}{logos/logo-uni-lu}

%%%% Fancy
\usepackage{fancyhdr}

%%%% Increase page length
\addtolength{\textheight}{1in}

%%%% Sections
\usepackage{sectsty}
\allsectionsfont{\sffamily}

%%%% Title
\newcommand{\titleone}{\textsf{Mutation Testing Survey:}}
\newcommand{\titletwo}{\textsf{FAQAS}}

\begin{document}

\pagestyle{fancy}
\renewcommand{\sectionmark}[1]{\markright{\textit{#1}}}

\renewcommand{\headrulewidth}{2pt}% 2pt header rule
\renewcommand{\headrule}{\hbox to\headwidth{%
  \color{sntblue}\leaders\hrule height \headrulewidth\hfill}}

\fancyhf{}

%\lhead{\fancyplain{}{\setlength{\unitlength}{1mm}
%\begin{picture}(0,0)
%\put(0,-4){\includegraphics[width=50pt]{logos/logo-snt}}
%\end{picture}}} 

\lhead{\fancyplain{}{\textit{Mutation Testing Survey}}}

\rhead{\fancyplain{}{\rightmark }}

\fancyfoot[C]{%
\begin{tikzpicture}[remember picture,overlay]
\path [fill=sntred]    ([xshift=88pt,yshift=20pt]current page.south west) rectangle
                       ([xshift=229pt,yshift=50pt] current page.south west);
\path [fill=sntpurple] ([xshift=229pt,yshift=20pt] current page.south west) rectangle
                       ([xshift=370pt,yshift=50pt] current page.south west);
\path [fill=sntblue]   ([xshift=370pt,yshift=20pt] current page.south west) rectangle
                       ([xshift=510pt,yshift=50pt] current page.south west);
\end{tikzpicture}
}

\fancyfoot[RO]{
\begin{tikzpicture}[remember picture,overlay]
\node [circle, ultra thick, fill=white, draw=sntblue] at ([xshift=530pt,yshift=35pt] current page.south west) {\thepage};
\end{tikzpicture}}

\fancyfoot[LE]{
\begin{tikzpicture}[remember picture,overlay]
\node [circle, ultra thick, fill=white, draw=sntblue] at ([xshift=65pt,yshift=35pt] current page.south west) {\thepage};
\end{tikzpicture}}

\thispagestyle{empty}

\begin{tikzpicture}[remember picture,overlay]
\path [fill=sntred]    ([xshift=30pt,yshift=20pt]current page.south west) rectangle
                       ([xshift=210pt,yshift=50pt] current page.south west);
\path [fill=sntpurple] ([xshift=210pt,yshift=20pt] current page.south west) rectangle
                       ([xshift=390pt,yshift=50pt] current page.south west);
\path [fill=sntblue]   ([xshift=390pt,yshift=20pt] current page.south west) rectangle
                       ([xshift=570pt,yshift=51pt] current page.south west);
\path [fill=unilublue] ([xshift=30pt,yshift= 50pt] current page.south west) --
                       ([xshift=570pt,yshift= 50pt] current page.south west)
                       [rounded corners=20pt] --
                       ([xshift=570pt,yshift=740pt] current page.south west)
                       [sharp corners] --
                       ([xshift=30pt,yshift=740pt] current page.south west);
\node [fill=white,rounded corners=0pt,inner xsep=6pt,inner ysep=3pt]
      at ([xshift=523pt,yshift=120pt] current page.south west)
      {\pgfuseimage{logo-uni-lu}};

\node [fill=white,rounded corners=2pt,inner xsep=6pt,inner ysep=3pt]
      at ([xshift=520pt,yshift=780pt] current page.south west)
      {\pgfuseimage{logo-snt}};

%\node [circle, fill=white, draw=sntblue] at ([xshift=550pt,yshift=35pt] current page.south west) {a};

\node[draw=none,fill=none,right] at (-1, -7){\color{white}\LARGE\bf\titleone};
\node[draw=none,fill=none,right] at (-1, -8){\color{white}\LARGE\bf\titletwo};
\node[draw=none,fill=none,right] at (-1, -9){\color{white}\LARGE\bf\titlethree};

\node[draw=none,fill=none,right] at (-1, -12){\color{white}\Large\textsf{O. Cornejo, F. Pastore, E. Viganò}};
\node[draw=none,fill=none,right] at (-1, -13){\color{white}\Large\textsf{Interdisciplinary Centre for Security, Reliability and Trust}};
\node[draw=none,fill=none,right] at (-1, -14){\color{white}\Large\textsf{University of Luxembourg}};
\node[draw=none,fill=none,right] at (11, -16){\color{white}\textsf{ITT-1-9873-ESA-FAQAS-SUTP}};
\node[draw=none,fill=none,right] at (11, -17){\color{white}\Large\textsf{Issue 4, Rev. 1}};
\node[draw=none,fill=none,right] at (11, -18){\color{white}\Large\textsf{\today}};
\node[draw=none,fill=none,right] at (-1, -24){\color{white}\tiny\textsf{EUROPEAN SPACE AGENCY. CONTRACT REPORT.}};
\node[draw=none,fill=none,right] at (-1, -24.3){\color{white}\tiny\textsf{The work described in this report was done under ESA contract. Responsibility for the contents resides in the author or organisation that prepared it.}};
\node[draw=none,fill=none,right] at (-1, -24.8){\color{white}\tiny\textsf{The copyright in this document is vested in the University of Luxembourg.}};
\node[draw=none,fill=none,right] at (-1, -25.1){\color{white}\tiny\textsf{This document may only be reproduced in whole or in part, or stored in a retrieval system,or transmitted in any form, or by any means electronic,}};
\node[draw=none,fill=none,right] at (-1, -25.4){\color{white}\tiny\textsf{mechanical, photocopying or otherwise, either with the prior permission of the University of Luxembourg or in accordance with the terms of ESTEC Contract No. 4000128969/19/NL/AS.}};

\node[draw=none,fill=none,right] at (-1, -24){\color{white}\tiny\textsf{EUROPEAN SPACE AGENCY. CONTRACT REPORT.}}; 
\node[draw=none,fill=none,right] at (-1, -24.3){\color{white}\tiny\textsf{The work described in this report was done under ESA contract. Responsibility for the contents resides in the author or organisation that prepared it.}};
\node[draw=none,fill=none,right] at (-1, -24.8){\color{white}\tiny\textsf{The copyright in this document is vested in the University of Luxembourg.}};
\node[draw=none,fill=none,right] at (-1, -25.1){\color{white}\tiny\textsf{This document may only be reproduced in whole or in part, or stored in a retrieval system,or transmitted in any form, or by any means electronic,}}; 
\node[draw=none,fill=none,right] at (-1, -25.4){\color{white}\tiny\textsf{mechanical, photocopying or otherwise, either with the prior permission of the University of Luxembourg or in accordance with the terms of ESTEC Contract No. 4000128969/19/NL/AS.}};


\end{tikzpicture}

\newpage


\section{Introduction}\label{sec:intro}

\cite{jia2010analysis}

\mbox{}\newpage

\newpage

\bibliographystyle{IEEEtran}
\bibliography{bibliography}


\end{document}  